\section{Semantic Retrieval and Text Similarity}

Ontologies provide structured, explicit representations of domain knowledge, 
but they depend on manually defined concepts and relations. 
This makes them well-suited for reasoning over known entities, 
yet limited in their ability to handle descriptions expressed in natural language 
or to retrieve knowledge that was not explicitly encoded. 
Semantic retrieval addresses this gap by operating on the meaning of text rather than on predefined structure. 
This section introduces the core concepts underlying this approach: how text can be represented as numerical vectors,
how similarity between texts can be measured in that representation, 
and how these techniques support ranking and matching against a corpus of scientific literature.

\subsection{Text Embeddings and Semantic Similarity}
\label{text-embeddings-semantic-similarity}

In vector semantics, text is represented as a numerical vector, 
a point in a high-dimensional space referred to as the embedding 
space \parencite[116]{jurafsky_martin2026speech_language_processing}. 
The position of a vector in this space reflects the semantic 
content of the text it represents. A central property of this 
representation is that semantically similar texts are mapped to 
nearby points in the embedding space, while unrelated texts are 
placed further apart. Figure~\ref{fig:embedding-space-example} illustrates 
this principle for a small set of example words.

\begin{figure}[H]
    \centering
    \includegraphics[width=0.9\textwidth]{Figures/02-background/embedding-space-example.png}
    \caption[Embedding space example]{A two-dimensional projection of example word embeddings.
    Semantically similar words cluster together in the embedding space,
    while unrelated words are positioned further apart.}
    \label{fig:embedding-space-example}
\end{figure}

To compare two vectors $\mathbf{u}$ and $\mathbf{v}$, cosine 
similarity is the most common metric used in natural language 
processing \parencite[Chapter~5.4]{jurafsky_martin2026speech_language_processing}. 
Rather than measuring absolute distance between two points, it 
measures the angle between two vectors, normalised for their length:

\begin{equation}
    \cos(\mathbf{u}, \mathbf{v}) = 
    \frac{\mathbf{u} \cdot \mathbf{v}}{\|\mathbf{u}\| \|\mathbf{v}\|}
    \label{eq:cosine-similarity}
\end{equation}

The result ranges from 1, indicating identical direction and 
maximum similarity, to 0 for orthogonal vectors with no similarity 
\parencite[Chapter~5.4]{jurafsky_martin2026speech_language_processing}. 
Normalising for vector length ensures that the similarity score 
reflects semantic content rather than vector magnitude. In 
practice, cosine distance, defined as $1 - \cos(\mathbf{u}, 
\mathbf{v})$, is often used when a distance metric is required, 
where lower values indicate greater similarity.

\subsection{Transformer-Based Sentence Embedding Models}
\label{sec:transformer-based-sentence-embedding-models}

Transformer-based language models represent text by processing 
input tokens through multiple layers of self-attention, producing 
contextual representations that capture relationships between 
words across the full input sequence 
\parencite{vaswani2023attentionneed}. Unlike earlier embedding 
approaches that assigned a fixed vector to each word regardless 
of context, transformer models produce representations that 
reflect the meaning of a word given its surrounding text.

However, standard transformer models such as BERT are not 
directly suited for computing sentence-level similarity. As 
noted by \textcite{reimers2019sentence_bert_sentence_embeddings_using_siamese_bert_networks}, 
comparing two sentences with BERT requires feeding both into 
the network simultaneously, which becomes computationally 
infeasible at scale. Finding the most similar pair in a 
collection of 10,000 sentences would require approximately 50 
million inference computations.

Sentence-BERT (SBERT) addresses this by fine-tuning BERT 
using a siamese network structure, producing fixed-size sentence 
embeddings that can be compared efficiently using cosine 
similarity \parencite{reimers2019sentence_bert_sentence_embeddings_using_siamese_bert_networks}. 
The key property is that semantically similar sentences are 
placed close together in the resulting embedding space, making 
SBERT suitable for tasks where embeddings for a large corpus 
can be precomputed and a new query compared against all of 
them in milliseconds.

\begin{figure}[H]
    \centering
    \includegraphics[width=0.7\textwidth]{Figures/02-background/SBERT-architecture.pdf}
    \caption[SBERT architecture]{SBERT architecture at inference time. Each sentence
    is encoded independently through a shared BERT network,
    producing vectors \textbf{u} and \textbf{v} that are compared
    using cosine similarity. Reproduced from
    \textcite{reimers2019sentence_bert_sentence_embeddings_using_siamese_bert_networks}
    under CC BY-SA 4.0.}
    \label{fig:sbert-architecture}
\end{figure}

The sentence-transformers library provides pretrained models 
based on this approach for encoding text into fixed-size 
vectors \parencite{sbert_docs}.

\subsection{Semantic Retrieval}

Information retrieval (IR) is the task of identifying and ranking 
documents from a collection according to their relevance to a 
user query \parencite[Chapter~11.1]{jurafsky_martin2026speech_language_processing}. 
Traditional approaches represent both documents and queries as 
sparse vectors of word counts, weighted by schemes such as tf-idf 
or BM25, and compute similarity based on shared terms 
\parencite[Chapter~11.1]{jurafsky_martin2026speech_language_processing}. While 
effective for many tasks, this approach treats documents as bags 
of words, meaning that word order and meaning are ignored, and 
only exact term matches contribute to the relevance score.

Dense retrieval addresses this limitation by representing 
documents and queries as embeddings computed from language 
models, rather than as word count vectors 
\parencite[Chapter~11.1]{jurafsky_martin2026speech_language_processing}. Because 
semantically similar texts are mapped to nearby points in the 
embedding space, a query can retrieve relevant documents even 
when they do not share exact terms with the query. This makes 
dense retrieval more flexible than keyword-based approaches for 
tasks where the user's information need is expressed in natural 
language rather than precise search terms.

Figure~\ref{fig:keyword-vs-dense} illustrates the distinction with a constructed example. 
The same query returns only Document~A under keyword-based retrieval, 
as it is the only document sharing exact terms with the query. 
Dense retrieval ranks all three documents by semantic similarity, 
including Documents~B and~C despite the absence of exact term overlap in these documents.

\begin{figure}[H]
    \centering
    \includegraphics[width=0.9\textwidth]{Figures/02-background/keyword-vs-dense.png}
    \caption[Keyword vs.\ dense retrieval]{Keyword-based versus dense retrieval for the same query.}
    \label{fig:keyword-vs-dense}
\end{figure}

\subsection{Cross-Encoder Reranking}
\label{sec:cross-encoder-reranking}

The sentence embedding approach described in 
Section~\ref{sec:transformer-based-sentence-embedding-models} is an example of a 
bi-encoder, where query and document are encoded independently 
into fixed-size vectors and similarity is measured using cosine 
similarity \parencite{reimers2019sentence_bert_sentence_embeddings_using_siamese_bert_networks}. 
Because the representations are computed independently, document 
embeddings can be precomputed and stored, making bi-encoders 
efficient for large-scale retrieval \parencite{thakur2021augmented_sbert_data_augmentation_method_improving_bi_encoders_pairwise_sentence_scoring_tasks}.

A cross-encoder takes a different approach, concatenating query 
and document into a single input and processing both jointly, 
which allows the model to apply attention across both texts 
simultaneously 
\parencite{reimers2019sentence_bert_sentence_embeddings_using_siamese_bert_networks}. 
This produces a direct relevance score rather than a similarity 
between independent vectors. Cross-encoders are generally more 
accurate than bi-encoders, but since no independent representations 
are computed, every query-document pair must be processed from 
scratch at query time, making cross-encoders impractical for 
retrieval over large document collections 
\parencite{thakur2021augmented_sbert_data_augmentation_method_improving_bi_encoders_pairwise_sentence_scoring_tasks}.

Figure~\ref{fig:bi-vs-cross-encoder} illustrates the architectural 
difference between the two approaches.

\begin{figure}[H]
    \centering
    \includegraphics[width=0.85\textwidth]{Figures/02-background/bi-vs-cross-encoder.png}
    \caption[Bi-encoder vs.\ cross-encoder]{Bi-encoder and cross-encoder architectures. The
    bi-encoder encodes query and document independently, producing
    vectors that are compared using cosine similarity. The
    cross-encoder concatenates both inputs and processes them
    jointly, enabling full attention between query and document
    inputs and producing a direct relevance score. Inspired by \textcite{kumar2024crossencoders}.}
    \label{fig:bi-vs-cross-encoder}
\end{figure}

A common approach to combining the strengths of both architectures 
is a two-stage retrieval pipeline 
\parencite{nogueira2019passage_re_ranking_with_bert}. In the first 
stage, a bi-encoder retrieves a set of candidate documents from 
the full corpus based on embedding similarity. In the second stage, 
a cross-encoder reranks these candidates by scoring each document 
against the query directly, producing a more precise relevance 
ranking \parencite{nogueira2019passage_re_ranking_with_bert}. This 
approach preserves the scalability of bi-encoder retrieval while 
benefiting from the higher accuracy of cross-encoder scoring.

\subsection{Scientific Literature as a Knowledge Source}

Scientific articles document how machine learning methods have 
been applied across a wide range of problems and contexts. In 
the domain of machine learning for engineering and manufacturing, 
a large number of such articles exist, covering diverse algorithms 
and applications \parencite{sala2018selecting_ml_algorithm}. 
This breadth of literature is itself a challenge. Relevant 
knowledge is spread across journals, conference proceedings, and 
research communities \parencite{adekoya2022applications_ai_em}, 
and the volume of available work makes it difficult for any 
individual to draw connections across topics and domains 
\parencite{olivetti2020data_driven_materials_research_enabled_natural_language_processing_information_extraction}. 
The result is that the literature creates extensive knowledge on 
the topic, while at the same time making
the selection of the right approach difficult
\parencite{sala2018selecting_ml_algorithm}.

Semantic retrieval offers a way to address this challenge by 
enabling systematic matching between a problem description and relevant articles. 
Article abstracts are a practical textual basis for this purpose. 
Text mining of scientific literature has mostly been carried out 
on abstracts due to their availability, and the fact that abstracts 
consist of shorter sentences presenting only the most 
important findings of a study 
\parencite{westergaard2018comprehensive_quantitative_comparison_text_mining_15_million_full_text_articles_versus_corresponding_abstracts}. 
While full-text articles contain richer information, they also 
include complex tables, references, and speculative discussion 
sections, making them less consistent as a basis for comparison 
across a large corpus 
\parencite{westergaard2018comprehensive_quantitative_comparison_text_mining_15_million_full_text_articles_versus_corresponding_abstracts}. 
Abstracts thus represent a structured and consistently available 
representation of each article's core content.

